\section{Conclusion}

This review set out to map the current state of AI-based supply chain risk
mitigation across three dimensions, the mechanisms through which AI is
applied, the risk categories it addresses, and the techniques it deploys,
drawing on 19 studies published between 2019 and 2026. What emerges is a
picture of a field that has grown considerably in both volume and technical
ambition over a relatively short period, shaped in large part by the
disruptions of the early 2020s and the data availability those disruptions
generated.

Across the included studies, early warning and forecasting stands out as
the most consistently pursued application of AI in supply chain risk
management, present in roughly 42\% of the corpus. Traditional machine
learning methods, Random Forest, Support Vector Machines, and gradient
boosting variants, account for the broadest deployment, owing to their
suitability for the structured tabular data that dominates operational
supply chain environments. Deep learning architectures, particularly LSTMs
and hybrid GNN-LSTM models, appear predominantly in studies dealing with
disruption and logistics risks, where temporal sequencing and network
topology carry analytical weight that simpler models cannot capture.
More recent contributions introduce large language models and multi-agent
architectures, reflecting a gradual expansion of the techniques considered
viable for operational risk contexts.

The review also surfaces a number of recurring limitations acknowledged
by the studies themselves. Coverage of cybersecurity, geopolitical, and
fraud risks remains thin relative to the attention given to disruption and
logistics risks. The empirical base is geographically concentrated, with
a substantial share of studies originating from Asian manufacturing and
logistics contexts, and with European and North American supply chain
environments largely absent from primary research. Most proposed models
are validated on historical batch datasets rather than live operational
data, a limitation that the authors of several included studies identify
as a barrier to real-world deployment. Interpretability, too, is widely
flagged as an unresolved challenge: despite the growing presence of
SHAP-based methods in recent contributions, the majority of reviewed
systems do not yet meet the transparency requirements of regulated or
safety-critical environments.

Future work signalled across the corpus points in several directions:
empirical validation in longitudinal, real-world settings; integration
with real-time data infrastructures such as IoT sensors and blockchain
audit trails; development of sector-specific frameworks for industries
currently underserved by the literature; and broader geographic coverage
to establish the cross-contextual generalisability of proposed approaches.
The field, as documented here, has made substantial progress in
demonstrating that AI can contribute meaningfully to supply chain risk
management; the work of establishing how, and under what conditions, that
contribution holds across diverse operational contexts remains largely
ahead of it.